% Part: applied-modal-logics
% Chapter: epistemic-logic
% Section: public-announcement-logic-lang

\documentclass[../../../include/open-logic-section]{subfiles}

\begin{document}

\olfileid{aml}{el}{pal}

\olsection{Logika Wara-Wara Umum}

Logika epistemik dinamis ngidini kita nggambarake carane kawruh para
agen owah manut wektu utawa nalika entuk informasi anyar. Akeh
logika iki nggambarake owah-owahan kawruh liwat \emph{prastawa}
informasi utawa \emph{panganyaran}. Jinis panganyaran kang paling
dhasar yaiku wara-wara umum: sawijining rumus diwara-warake kanthi
bener lan kabeh agen nyekseni prastawa iku bebarengan. Kanggo iki,
kita ngambakake basane kaya mangkene.

\begin{defn}
Upamane $G$ himpunan lambang agen. Basa dhasar logika epistemik
multiagen kanthi wara-wara umum ngemot
\begin{enumerate}
  \tagitem{prvFalse}{Konstanta proposisional kanggo !!{falsity}~$\lfalse$.}{}
  \tagitem{prvTrue}{Konstanta proposisional kanggo !!{truth}~$\ltrue$.}{}
  \item Sawijining himpunan kang !!{denumerable}s saka !!{propositional variable}s: $\Obj
    p_0$, $\Obj p_1$, $\Obj p_2$, \dots
  \item Konektif proposisional: \startycommalist
  \iftag{prvNot}{\ycomma $\lnot$ (negasi)}{}%
  \iftag{prvAnd}{\ycomma $\land$ (konjungsi)}{}%
  \iftag{prvOr}{\ycomma $\lor$ (disjungsi)}{}%
  \iftag{prvIf}{\ycomma $\lif$ (!!{conditional})}{}%
  \iftag{prvIff}{\ycomma $\liff$ (!!{biconditional})}{}.
  \item Operator kawruh $\Knows_a$ kang $a \in G$.
  \item Operator wara-wara umum $[!B]$ kang $!B$ iku !!a{formula}.
\end{enumerate}
\end{defn}

Operator wara-wara umum tumindak minangka operator kothak, mula
definisi basa kanthi induksi diwenehake kaya mangkene:

\begin{defn}
  \emph{!!^{formula}s} ing basa epistemik ditetepake kanthi induksi
  kaya mangkene:
  \begin{enumerate}
  \tagitem{prvFalse}{$\lfalse$ iku !!{formula} atomik.}{}

  \tagitem{prvTrue}{$\ltrue$ iku !!{formula} atomik.}{}

  \item Saben variabel proposisional $\Obj p_i$ iku !!{formula} (atomik).

  \tagitem{prvNot}{Yen $!A$ iku !!a{formula}, $\lnot !A$ uga
  !!a{formula}.}{}

  \tagitem{prvAnd}{Yen $!A$ lan $!B$ iku !!{formula}s, $(!A \land
  !B)$ uga !!a{formula}.}{}

  \tagitem{prvOr}{Yen $!A$ lan $!B$ iku !!{formula}s, $(!A \lor !B)$
  uga !!a{formula}.}{}

  \tagitem{prvIf}{Yen $!A$ lan $!B$ iku !!{formula}s, $(!A \lif !B)$
  uga !!a{formula}.}{}

  \tagitem{prvIff}{Yen $!A$ lan $!B$ iku !!{formula}s, $(!A \liff !B)$
  uga !!a{formula}.}{}

  \item Yen $!A$ iku !!a{formula} lan $a \in G$, $\Knows_a !A$
  uga !!a{formula}.
  
  \item Yen $!A$ lan $!B$ iku !!{formula}s, $[!A] !B$ uga !!a{formula}.

  \tagitem{limitClause}{Ora ana liyane kang dadi !!a{formula}.}{}
\end{enumerate}
\end{defn}

!!{formula} $[!A] !B$ dimaksudake diwaca ``Sawise $!A$ diwara-warake
kanthi bener, $!B$~bener.'' Kadhangkala uga migunani kanggo ngrembug
kawruh umum ing konteks wara-wara umum, mula basa iki uga bisa ngemot
operator kawruh umum~$\CKnows_G
!A$.

\end{document}
